Os dados de referência mostram que as densidades dos três fluidos são todas da mesma ordem de grandeza, enquanto suas viscosidades diferem dramaticamente.

\begin{itemize}
    \item \textbf{Água:} $\rho\approx1{,}00\times10^{3}\ \text{kg/m}^{3}$ e $\mu\approx1{,}0\times10^{-3}\ \text{Pa}\cdot\text{s}$.
    \item \textbf{Mel:} $\rho\approx1{,}40\times10^{3}\ \text{kg/m}^{3}$ e $\mu$ tipicamente entre $2$ e $10\ \text{Pa}\cdot\text{s}$.
    \item \textbf{Manteiga de amendoim:} $\rho$ varia aproximadamente entre $1{,}00\times10^{3}$ e $1{,}30\times10^{3}\ \text{kg/m}^{3}$; a viscosidade pode chegar a centenas de pascal-segundo, frequentemente $\gtrsim250\ \text{Pa}\cdot\text{s}$ dependendo da temperatura e do tipo.
\end{itemize}

\textbf{Comparação:}
\begin{itemize}
    \item A densidade da água, do mel e da manteiga de amendoim difere em menos de um fator de 2.
    \item A viscosidade difere em vários milhares de vezes: a água tem $10^{-3}\ \text{Pa}\cdot\text{s}$, o mel tem alguns pascal-segundo e a manteiga de amendoim pode exceder $250\ \text{Pa}\cdot\text{s}$.
\end{itemize}

Isso confirma que, na modelagem do misturador, a viscosidade é o parâmetro dominante para explicar por que a manteiga de amendoim requer muito mais força do que a água, enquanto a densidade por si só não explica essa diferença extrema.
